\documentclass[10pt,a4paper]{amsart}
\usepackage{amsmath,amssymb,amsthm,graphicx}
% For commutative diagrams we use Xy-pic
\usepackage[all]{xy}

% We use 2cell for 2-commutative diagrams.
\xyoption{2cell}
% Theorem environments.
%
\theoremstyle{plain}
\newtheorem{theorem}[subsection]{Theorem}
\newtheorem{proposition}[subsection]{Proposition}
\newtheorem{lemma}[subsection]{Lemma}

\theoremstyle{definition}
\newtheorem{definition}[subsection]{Definition}
\newtheorem{example}[subsection]{Example}
\newtheorem{exercise}[subsection]{Exercise}
\newtheorem{situation}[subsection]{Situation}

\theoremstyle{remark}
\newtheorem{remark}[subsection]{Remark}
\newtheorem{remarks}[subsection]{Remarks}

\numberwithin{equation}{subsection}

% Macros
%
\def\lim{\mathop{\mathrm{lim}}\nolimits}
\def\colim{\mathop{\mathrm{colim}}\nolimits}
\def\Spec{\mathop{\mathrm{Spec}}}
\def\Hom{\mathop{\mathrm{Hom}}\nolimits}
\def\Ext{\mathop{\mathrm{Ext}}\nolimits}
\def\SheafHom{\mathop{\mathcal{H}\!\mathit{om}}\nolimits}
\def\SheafExt{\mathop{\mathcal{E}\!\mathit{xt}}\nolimits}
\def\Sch{\mathit{Sch}}
\def\Mor{\mathop{\mathrm{Mor}}\nolimits}
\def\Ob{\mathop{\mathrm{Ob}}\nolimits}
\def\Sh{\mathop{\mathit{Sh}}\nolimits}
\def\NL{\mathop{N\!L}\nolimits}
\def\CH{\mathop{\mathrm{CH}}\nolimits}
\def\proetale{{pro\text{-}\acute{e}tale}}
\def\etale{{\acute{e}tale}}
\def\QCoh{\mathit{QCoh}}
\def\Ker{\mathop{\mathrm{Ker}}}
\def\Im{\mathop{\mathrm{Im}}}
\def\Coker{\mathop{\mathrm{Coker}}}
\def\Coim{\mathop{\mathrm{Coim}}}

% Boxtimes
%
\DeclareMathSymbol{\boxtimes}{\mathbin}{AMSa}{"02}

%
% Macros for moduli stacks/spaces
%
\def\QCohstack{\mathcal{QC}\!\mathit{oh}}
\def\Cohstack{\mathcal{C}\!\mathit{oh}}
\def\Spacesstack{\mathcal{S}\!\mathit{paces}}
\def\Quotfunctor{\mathrm{Quot}}
\def\Hilbfunctor{\mathrm{Hilb}}
\def\Curvesstack{\mathcal{C}\!\mathit{urves}}
\def\Polarizedstack{\mathcal{P}\!\mathit{olarized}}
\def\Complexesstack{\mathcal{C}\!\mathit{omplexes}}
% \Pic is the operator that assigns to X its picard group, usage \Pic(X)
% \Picardstack_{X/B} denotes the Picard stack of X over B
% \Picardfunctor_{X/B} denotes the Picard functor of X over B
\def\Pic{\mathop{\mathrm{Pic}}\nolimits}
\def\Picardstack{\mathcal{P}\!\mathit{ic}}
\def\Picardfunctor{\mathrm{Pic}}
\def\Deformationcategory{\mathcal{D}\!\mathit{ef}}
\newenvironment{reference}{}{}
\newenvironment{slogan}{}{}
\newenvironment{history}{}{}
\usepackage{cleveref}
\crefname{originalequation}{original equation}{original equations}
\Crefname{originalequation}{Original equation}{Original equations}
\makeatletter
\newlabel{equation-invariants}{{03BF}{source}{}{sourcepointer-1}{}}
\newlabel{equation-invariants@cref}{{[originalequation][0][]03BF}{[1][0][]source}{}{}{}}
\makeatother
\crefname{originallemma}{original lemma}{original lemmas}
\Crefname{originallemma}{Original lemma}{Original lemmas}
\makeatletter
\newlabel{lemma-finite-locally-free-disjoint-free}{{03BI}{source}{}{sourcepointer-2}{}}
\newlabel{lemma-finite-locally-free-disjoint-free@cref}{{[originallemma][0][]03BI}{[1][0][]source}{}{}{}}
\makeatother
\crefname{originallemma}{original lemma}{original lemmas}
\Crefname{originallemma}{Original lemma}{Original lemmas}
\makeatletter
\newlabel{lemma-determinant-trick}{{03BH}{source}{}{sourcepointer-3}{}}
\newlabel{lemma-determinant-trick@cref}{{[originallemma][0][]03BH}{[1][0][]source}{}{}{}}
\makeatother
\crefname{originallemma}{original lemma}{original lemmas}
\Crefname{originallemma}{Original lemma}{Original lemmas}
\makeatletter
\newlabel{algebra-lemma-charpoly}{{00DX}{source}{}{sourcepointer-4}{}}
\newlabel{algebra-lemma-charpoly@cref}{{[originallemma][0][]00DX}{[1][0][]source}{}{}{}}
\makeatother
\crefname{originallemma}{original lemma}{original lemmas}
\Crefname{originallemma}{Original lemma}{Original lemmas}
\makeatletter
\newlabel{algebra-lemma-integral-overring-surjective}{{00GQ}{source}{}{sourcepointer-5}{}}
\newlabel{algebra-lemma-integral-overring-surjective@cref}{{[originallemma][0][]00GQ}{[1][0][]source}{}{}{}}
\makeatother
\crefname{originallemma}{original lemma}{original lemmas}
\Crefname{originallemma}{Original lemma}{Original lemmas}
\makeatletter
\newlabel{morphisms-lemma-base-change-surjective}{{01S1}{source}{}{sourcepointer-6}{}}
\newlabel{morphisms-lemma-base-change-surjective@cref}{{[originallemma][0][]01S1}{[1][0][]source}{}{}{}}
\makeatother
\crefname{originallemma}{original lemma}{original lemmas}
\Crefname{originallemma}{Original lemma}{Original lemmas}
\makeatletter
\newlabel{algebra-lemma-image-dense-generic-points}{{00FL}{source}{}{sourcepointer-7}{}}
\newlabel{algebra-lemma-image-dense-generic-points@cref}{{[originallemma][0][]00FL}{[1][0][]source}{}{}{}}
\makeatother
\crefname{originallemma}{original lemma}{original lemmas}
\Crefname{originallemma}{Original lemma}{Original lemmas}
\makeatletter
\newlabel{algebra-lemma-silly}{{00DS}{source}{}{sourcepointer-8}{}}
\newlabel{algebra-lemma-silly@cref}{{[originallemma][0][]00DS}{[1][0][]source}{}{}{}}
\makeatother
\crefname{originallemma}{original lemma}{original lemmas}
\Crefname{originallemma}{Original lemma}{Original lemmas}
\makeatletter
\newlabel{limits-lemma-descend-opens}{{01Z4}{source}{}{sourcepointer-9}{}}
\newlabel{limits-lemma-descend-opens@cref}{{[originallemma][0][]01Z4}{[1][0][]source}{}{}{}}
\makeatother
\crefname{originallemma}{original lemma}{original lemmas}
\Crefname{originallemma}{Original lemma}{Original lemmas}
\makeatletter
\newlabel{limits-lemma-descend-isomorphism}{{081E}{source}{}{sourcepointer-10}{}}
\newlabel{limits-lemma-descend-isomorphism@cref}{{[originallemma][0][]081E}{[1][0][]source}{}{}{}}
\makeatother
\crefname{originallemma}{original lemma}{original lemmas}
\Crefname{originallemma}{Original lemma}{Original lemmas}
\makeatletter
\newlabel{morphisms-lemma-set-points-where-fibres-etale}{{0476}{source}{}{sourcepointer-11}{}}
\newlabel{morphisms-lemma-set-points-where-fibres-etale@cref}{{[originallemma][0][]0476}{[1][0][]source}{}{}{}}
\makeatother
\crefname{originallemma}{original lemma}{original lemmas}
\Crefname{originallemma}{Original lemma}{Original lemmas}
\makeatletter
\newlabel{algebra-lemma-finite-residue-extension-closed}{{00GA}{source}{}{sourcepointer-12}{}}
\newlabel{algebra-lemma-finite-residue-extension-closed@cref}{{[originallemma][0][]00GA}{[1][0][]source}{}{}{}}
\makeatother
\crefname{originallemma}{original lemma}{original lemmas}
\Crefname{originallemma}{Original lemma}{Original lemmas}
\makeatletter
\newlabel{more-algebra-lemma-integral-over-henselian-pair}{{09XK}{source}{}{sourcepointer-13}{}}
\newlabel{more-algebra-lemma-integral-over-henselian-pair@cref}{{[originallemma][0][]09XK}{[1][0][]source}{}{}{}}
\makeatother
\crefname{originallemma}{original lemma}{original lemmas}
\Crefname{originallemma}{Original lemma}{Original lemmas}
\makeatletter
\newlabel{more-algebra-lemma-change-ideal-henselian-pair}{{09XJ}{source}{}{sourcepointer-14}{}}
\newlabel{more-algebra-lemma-change-ideal-henselian-pair@cref}{{[originallemma][0][]09XJ}{[1][0][]source}{}{}{}}
\makeatother
\crefname{originallemma}{original lemma}{original lemmas}
\Crefname{originallemma}{Original lemma}{Original lemmas}
\makeatletter
\newlabel{more-algebra-lemma-characterize-henselian-pair}{{09XI}{source}{}{sourcepointer-15}{}}
\newlabel{more-algebra-lemma-characterize-henselian-pair@cref}{{[originallemma][0][]09XI}{[1][0][]source}{}{}{}}
\makeatother
\crefname{originallemma}{original lemma}{original lemmas}
\Crefname{originallemma}{Original lemma}{Original lemmas}
\makeatletter
\newlabel{morphisms-lemma-etale-permanence}{{02GW}{source}{}{sourcepointer-16}{}}
\newlabel{morphisms-lemma-etale-permanence@cref}{{[originallemma][0][]02GW}{[1][0][]source}{}{}{}}
\makeatother
\crefname{originallemma}{original lemma}{original lemmas}
\Crefname{originallemma}{Original lemma}{Original lemmas}
\makeatletter
\newlabel{descent-lemma-universally-injective-etale-open-immersion}{{02LC}{source}{}{sourcepointer-17}{}}
\newlabel{descent-lemma-universally-injective-etale-open-immersion@cref}{{[originallemma][0][]02LC}{[1][0][]source}{}{}{}}
\makeatother
\title{Stabilizer-preserving etale maps of finite-flat groupoids descend to etale invariant quotients}
\author{The Stacks Project Authors}
\date{stacks-project:groupoids.tex:local-snapshot-20261011; GNU FDL 1.2 or later}
\begin{document}
\maketitle
\noindent\textit{Editorial note.} Public299 treats inertia invariants of a finite abstract group acting on a ring. Here two finite-flat affine groupoids and a stabilizer-preserving morphism are given, and the conclusion includes a Cartesian atlas square.\par\medskip
\noindent\textit{Editorial note.} Public1742 constructs an fppf invariant quotient for an equivalence relation; it does not establish this etale change theorem with stabilizers.\par\medskip
\noindent\textit{Allowed standard prerequisites (editorial).}\begin{itemize}
\item Explicit prerequisite boundary. Affine groupoid schemes, finite locally free maps, strict henselization and etale descent, finite invariant-ring quotients and stabilizers. The source-local integrality, invariant base change and point-orbit arguments are retained.
\end{itemize}\medskip
\ifdefined\hypertarget\hypertarget{sourcepointer-1}{}\fi
\noindent\textit{Source reference.} Original equation 03BF: focused-seeds/AGP/sources/stacks-snapshot-20261011/tags.txt, source lines 4346--4346. Stable Stacks tag for an expressly allowed foundational or notational input; the counted construction remains the two concrete core obligations. The retained local core is not replaced by this pointer. See standard\_prerequisites.\par\medskip
\ifdefined\hypertarget\hypertarget{sourcepointer-2}{}\fi
\noindent\textit{Source reference.} Original lemma 03BI: focused-seeds/AGP/sources/stacks-snapshot-20261011/tags.txt, source lines 4349--4349. Stable Stacks tag for an expressly allowed foundational or notational input; the counted construction remains the two concrete core obligations. The retained local core is not replaced by this pointer. See standard\_prerequisites.\par\medskip
\ifdefined\hypertarget\hypertarget{sourcepointer-3}{}\fi
\noindent\textit{Source reference.} Original lemma 03BH: focused-seeds/AGP/sources/stacks-snapshot-20261011/tags.txt, source lines 4348--4348. Stable Stacks tag for an expressly allowed foundational or notational input; the counted construction remains the two concrete core obligations. The retained local core is not replaced by this pointer. See standard\_prerequisites.\par\medskip
\ifdefined\hypertarget\hypertarget{sourcepointer-4}{}\fi
\noindent\textit{Source reference.} Original lemma 00DX: focused-seeds/AGP/sources/stacks-snapshot-20261011/tags.txt, source lines 514--514. Stable Stacks tag for an expressly allowed foundational or notational input; the counted construction remains the two concrete core obligations. The retained local core is not replaced by this pointer. See standard\_prerequisites.\par\medskip
\ifdefined\hypertarget\hypertarget{sourcepointer-5}{}\fi
\noindent\textit{Source reference.} Original lemma 00GQ: focused-seeds/AGP/sources/stacks-snapshot-20261011/tags.txt, source lines 616--616. Stable Stacks tag for an expressly allowed foundational or notational input; the counted construction remains the two concrete core obligations. The retained local core is not replaced by this pointer. See standard\_prerequisites.\par\medskip
\ifdefined\hypertarget\hypertarget{sourcepointer-6}{}\fi
\noindent\textit{Source reference.} Original lemma 01S1: focused-seeds/AGP/sources/stacks-snapshot-20261011/tags.txt, source lines 2322--2322. Stable Stacks tag for an expressly allowed foundational or notational input; the counted construction remains the two concrete core obligations. The retained local core is not replaced by this pointer. See standard\_prerequisites.\par\medskip
\ifdefined\hypertarget\hypertarget{sourcepointer-7}{}\fi
\noindent\textit{Source reference.} Original lemma 00FL: focused-seeds/AGP/sources/stacks-snapshot-20261011/tags.txt, source lines 575--575. Stable Stacks tag for an expressly allowed foundational or notational input; the counted construction remains the two concrete core obligations. The retained local core is not replaced by this pointer. See standard\_prerequisites.\par\medskip
\ifdefined\hypertarget\hypertarget{sourcepointer-8}{}\fi
\noindent\textit{Source reference.} Original lemma 00DS: focused-seeds/AGP/sources/stacks-snapshot-20261011/tags.txt, source lines 508--508. Stable Stacks tag for an expressly allowed foundational or notational input; the counted construction remains the two concrete core obligations. The retained local core is not replaced by this pointer. See standard\_prerequisites.\par\medskip
\ifdefined\hypertarget\hypertarget{sourcepointer-9}{}\fi
\noindent\textit{Source reference.} Original lemma 01Z4: focused-seeds/AGP/sources/stacks-snapshot-20261011/tags.txt, source lines 2579--2579. Stable Stacks tag for an expressly allowed foundational or notational input; the counted construction remains the two concrete core obligations. The retained local core is not replaced by this pointer. See standard\_prerequisites.\par\medskip
\ifdefined\hypertarget\hypertarget{sourcepointer-10}{}\fi
\noindent\textit{Source reference.} Original lemma 081E: focused-seeds/AGP/sources/stacks-snapshot-20261011/tags.txt, source lines 10214--10214. Stable Stacks tag for an expressly allowed foundational or notational input; the counted construction remains the two concrete core obligations. The retained local core is not replaced by this pointer. See standard\_prerequisites.\par\medskip
\ifdefined\hypertarget\hypertarget{sourcepointer-11}{}\fi
\noindent\textit{Source reference.} Original lemma 0476: focused-seeds/AGP/sources/stacks-snapshot-20261011/tags.txt, source lines 5497--5497. Stable Stacks tag for an expressly allowed foundational or notational input; the counted construction remains the two concrete core obligations. The retained local core is not replaced by this pointer. See standard\_prerequisites.\par\medskip
\ifdefined\hypertarget\hypertarget{sourcepointer-12}{}\fi
\noindent\textit{Source reference.} Original lemma 00GA: focused-seeds/AGP/sources/stacks-snapshot-20261011/tags.txt, source lines 600--600. Stable Stacks tag for an expressly allowed foundational or notational input; the counted construction remains the two concrete core obligations. The retained local core is not replaced by this pointer. See standard\_prerequisites.\par\medskip
\ifdefined\hypertarget\hypertarget{sourcepointer-13}{}\fi
\noindent\textit{Source reference.} Original lemma 09XK: focused-seeds/AGP/sources/stacks-snapshot-20261011/tags.txt, source lines 12533--12533. Stable Stacks tag for an expressly allowed foundational or notational input; the counted construction remains the two concrete core obligations. The retained local core is not replaced by this pointer. See standard\_prerequisites.\par\medskip
\ifdefined\hypertarget\hypertarget{sourcepointer-14}{}\fi
\noindent\textit{Source reference.} Original lemma 09XJ: focused-seeds/AGP/sources/stacks-snapshot-20261011/tags.txt, source lines 12532--12532. Stable Stacks tag for an expressly allowed foundational or notational input; the counted construction remains the two concrete core obligations. The retained local core is not replaced by this pointer. See standard\_prerequisites.\par\medskip
\ifdefined\hypertarget\hypertarget{sourcepointer-15}{}\fi
\noindent\textit{Source reference.} Original lemma 09XI: focused-seeds/AGP/sources/stacks-snapshot-20261011/tags.txt, source lines 12531--12531. Stable Stacks tag for an expressly allowed foundational or notational input; the counted construction remains the two concrete core obligations. The retained local core is not replaced by this pointer. See standard\_prerequisites.\par\medskip
\ifdefined\hypertarget\hypertarget{sourcepointer-16}{}\fi
\noindent\textit{Source reference.} Original lemma 02GW: focused-seeds/AGP/sources/stacks-snapshot-20261011/tags.txt, source lines 3243--3243. Stable Stacks tag for an expressly allowed foundational or notational input; the counted construction remains the two concrete core obligations. The retained local core is not replaced by this pointer. See standard\_prerequisites.\par\medskip
\ifdefined\hypertarget\hypertarget{sourcepointer-17}{}\fi
\noindent\textit{Source reference.} Original lemma 02LC: focused-seeds/AGP/sources/stacks-snapshot-20261011/tags.txt, source lines 3403--3403. Stable Stacks tag for an expressly allowed foundational or notational input; the counted construction remains the two concrete core obligations. The retained local core is not replaced by this pointer. See standard\_prerequisites.\par\medskip
\begin{lemma}
\label{lemma-etale}
Let $S$ be a scheme. Let $f : (U', R', s', t') \to (U, R, s, t, c)$ be a
morphism of groupoid schemes over $S$.
\begin{enumerate}
\item $U$, $R$, $U'$, $R'$ are affine,
\item $s, t, s', t'$ are finite locally free,
\item the diagrams
$$
\xymatrix{
R' \ar[d]_{s'} \ar[r]_f & R \ar[d]^s \\
U' \ar[r]^f & U
}
\quad
\quad
\xymatrix{
R' \ar[d]_{t'} \ar[r]_f & R \ar[d]^t \\
U' \ar[r]^f & U
}
\quad
\quad
\xymatrix{
G' \ar[d] \ar[r]_f & G \ar[d] \\
U' \ar[r]^f & U
}
$$
are cartesian where $G$ and $G'$ are the stabilizer group schemes, and
\item $f : U' \to U$ is \'etale.
\end{enumerate}
Then the map $C \to C'$ from the $R$-invariant functions on $U$
to the $R'$-invariant functions on $U'$ is \'etale and
$U' = \Spec(C') \times_{\Spec(C)} U$.
\end{lemma}

\begin{lemma}
\label{lemma-integral-over-invariants}
Let $S$ be a scheme. Let $(U, R, s, t, c)$ be a groupoid scheme over $S$.
Assume $U = \Spec(A)$ and $R = \Spec(B)$ are affine and
$s, t : R \to U$ finite locally free.
Let $C \subset A$ be as in (\ref{equation-invariants}).
Then $A$ is integral over $C$.
\end{lemma}

\begin{proof}
First, by Lemma \ref{lemma-finite-locally-free-disjoint-free}
we know that $(U, R, s, t, c)$ is a disjoint union
of groupoid schemes $(U_r, R_r, s, t, c)$ such that each $s, t : R_r \to U_r$
has constant rank $r$. As $U$ is quasi-compact, we have $U_r = \emptyset$ for
almost all $r$. It suffices to prove the lemma for each $(U_r, R_r, s, t, c)$
and hence we may assume that $s, t$ are finite locally free of rank $r$.

\medskip\noindent
Assume that $s, t$ are finite locally free of rank $r$.
Let $f \in A$. Consider the element $x - f \in A[x]$, where we think
of $x$ as the coordinate on $\mathbf{A}^1$.
Since
$$
(U \times \mathbf{A}^1, R \times \mathbf{A}^1,
s \times \text{id}_{\mathbf{A}^1},
t \times \text{id}_{\mathbf{A}^1},
c \times \text{id}_{\mathbf{A}^1})
$$
is also a groupoid scheme with finite source and target, we may apply
Lemma \ref{lemma-determinant-trick} to it and we see that
$P(x) = \text{Norm}_{s^\sharp}(t^\sharp(x - f))$
is an element of $C[x]$. Because $s^\sharp : A \to B$ is finite locally
free of rank $r$ we see that $P$ is monic of degree $r$.
Moreover $P(f) = 0$ by Cayley-Hamilton
(Algebra, Lemma \ref{algebra-lemma-charpoly}).
More precisely, Cayley-Hamilton says the polynomial
$s^\sharp(P) \in B[x]$ annihilates the element $t^\sharp(f)$.
Since the coefficients of $P$ are in $C$ we get that
$t^\sharp(P(f)) = 0$ in $B$. Since $t$ is faithfully flat,
we get $P(f) = 0$.
\end{proof}

\begin{lemma}
\label{lemma-invariants-base-change}
Let $S$ be a scheme. Let $(U, R, s, t, c)$ be a groupoid scheme over $S$.
Assume $U = \Spec(A)$ and $R = \Spec(B)$ are affine and
$s, t : R \to U$ finite locally free. Let $C \subset A$ be as in
(\ref{equation-invariants}). Let $C \to C'$ be a ring map, and set
$U' = \Spec(A \otimes_C C')$,
$R' = \Spec(B \otimes_C C')$.
Then
\begin{enumerate}
\item The maps $s, t, c$ induce maps $s', t', c'$ such that
$(U', R', s', t', c')$ is a groupoid scheme. Let $C^1 \subset A'$
be the $R'$-invariant functions on $U'$.
\item The canonical map $\varphi : C' \to C^1$ satisfies
\begin{enumerate}
\item for every $f \in C^1$ there exists an $n > 0$ and a
polynomial $P \in C'[x]$ whose image in $C^1[x]$ is
$(x - f)^n$, and
\item for every $f \in \Ker(\varphi)$ there exists
an $n > 0$ such that $f^n = 0$.
\end{enumerate}
\item If $C \to C'$ is flat then $\varphi$ is an isomorphism.
\end{enumerate}
\end{lemma}

\begin{proof}
The proof of part (1) is omitted. Let us denote $A' = A \otimes_C C'$ and
$B' = B \otimes_C C'$. Then we have
$$
C^1
= \{a \in A' \mid (t')^\sharp(a) = (s')^\sharp(a) \}
= \{a \in A \otimes_C C' \mid t^\sharp \otimes 1(a) = s^\sharp \otimes 1(a) \}.
$$
In other words, $C^1$ is the kernel of the difference map
$(t^\sharp - s^\sharp) \otimes 1$ which is just the base change
of the $C$-linear map $t^\sharp - s^\sharp : A \to B$ by $C \to C'$.
Hence (3) follows.

\medskip\noindent
Proof of part (2)(b). Since $C \to A$ is integral
(Lemma \ref{lemma-integral-over-invariants}) and injective we see that
$\Spec(A) \to \Spec(C)$ is surjective, see
Algebra, Lemma \ref{algebra-lemma-integral-overring-surjective}.
Thus also $\Spec(A') \to \Spec(C')$ is surjective
as a base change of a surjective morphism
(Morphisms, Lemma \ref{morphisms-lemma-base-change-surjective}).
Hence $\Spec(C^1) \to \Spec(C')$ is surjective also.
This implies (2)(b) holds for example by
Algebra, Lemma \ref{algebra-lemma-image-dense-generic-points}.

\medskip\noindent
Proof of part (2)(a). By Lemma \ref{lemma-finite-locally-free-disjoint-free}
our groupoid scheme $(U, R, s, t, c)$ decomposes as a finite disjoint union
of groupoid schemes $(U_r, R_r, s, t, c)$ such that $s, t : R_r \to U_r$
are finite locally free of rank $r$. Pulling back by $U' = \Spec(C') \to U$
we obtain a similar decomposition of $U'$ and $U^1 = \Spec(C^1)$.
We will show in the next paragraph that (2)(a) holds for the corresponding
system of rings $A_r, B_r, C_r, C'_r, C^1_r$ with $n = r$.
Then given $f \in C^1$ let $P_r \in C_r[x]$ be the polynomial
whose image in $C^1_r[x]$ is the image of $(x - f)^r$.
Choosing a sufficiently divisible integer $n$ we see that
there is a polynomial $P \in C'[x]$ whose image in $C^1[x]$ is
$(x - f)^n$; namely, we take $P$ to be the unique element of
$C'[x]$ whose image in $C'_r[x]$ is $P_r^{n/r}$.

\medskip\noindent
In this paragraph we prove (2)(a) in case the ring maps
$s^\sharp, t^\sharp : A \to B$ are finite locally free of a fixed rank $r$.
Let $f \in C^1 \subset A' = A \otimes_C C'$. Choose a flat
$C$-algebra $D$ and a surjection $D \to C'$. Choose a lift
$g \in A \otimes_C D$ of $f$.
Consider the polynomial
$$
P = \text{Norm}_{s^\sharp \otimes 1}(t^\sharp \otimes 1(x - g))
$$
in $(A \otimes_C D)[x]$. By Lemma \ref{lemma-determinant-trick}
and part (3) of the current lemma the coefficients of $P$ are in $D$
(compare with the proof of Lemma \ref{lemma-integral-over-invariants}).
On the other hand, the image of $P$ in $(A \otimes_C C')[x]$ is
$(x - f)^r$ because $t^\sharp \otimes 1(x - f) = s^\sharp(x - f)$
and $s^\sharp$ is finite locally free of rank $r$.
This proves what we want with $P$ as in the statement (2)(a)
given by the image of our $P$ under the map $D[x] \to C'[x]$.
\end{proof}

\begin{lemma}
\label{lemma-points}
Let $S$ be a scheme. Let $(U, R, s, t, c)$ be a groupoid scheme over $S$.
Assume $U = \Spec(A)$ and $R = \Spec(B)$ are affine and
$s, t : R \to U$ finite locally free. Let $C \subset A$ be as in
(\ref{equation-invariants}). Then $U \to M = \Spec(C)$ has
the following properties:
\begin{enumerate}
\item the map on points $|U| \to |M|$ is surjective and
$u_0, u_1 \in |U|$ map to the same point if and only if
there exists a $r \in |R|$ with $t(r) = u_0$ and $s(r) = u_1$, in
a formula
$$
|M| = |U|/|R|
$$
\item for any algebraically closed field $k$ we have
$$
M(k) = U(k)/R(k)
$$
\end{enumerate}
\end{lemma}

\begin{proof}
Since $C \to A$ is integral (Lemma \ref{lemma-integral-over-invariants})
and injective we see that $\Spec(A) \to \Spec(C)$ is surjective, see
Algebra, Lemma \ref{algebra-lemma-integral-overring-surjective}.
Thus $|U| \to |M|$ is surjective.

\medskip\noindent
Let $k$ be an algebraically closed field and let $C \to k$ be a ring map.
Since surjective morphisms are preserved under base change
(Morphisms, Lemma \ref{morphisms-lemma-base-change-surjective})
we see that $A \otimes_C k$ is not zero. Now $k \subset A \otimes_C k$ is a
nonzero integral extension. Hence any residue field of $A \otimes_C k$
is an algebraic extension of $k$, hence equal to $k$. Thus we see that
$U(k) \to M(k)$ is surjective.

\medskip\noindent
Let $a_0, a_1 : A \to k$ be two ring maps. If there exists a ring map
$b : B \to k$ such that $a_0 = b \circ t^\sharp$ and $a_1 = b \circ s^\sharp$
then we see that $a_0|_C = a_1|_C$ by definition. Thus the map
$U(k) \to M(k)$ equalizes the two maps $R(k) \to U(k)$.
Conversely, suppose that $a_0|_C = a_1|_C$. Let us name this algebra
map $c : C \to k$. Consider the diagram
$$
\xymatrix{
& &
B \ar@{-->}[lld] \\
k & &
A
\ar@<0.5ex>[ll]^{a_0}
\ar@<-0.5ex>[ll]_{a_1}
\ar@<1ex>[u]
\ar@<-1ex>[u] \\
& &
C \ar[u] \ar[llu]^c
}
$$
If we can construct a dotted arrow making the diagram commute, then
the proof of part (2) of the lemma is complete. Since $s : A \to B$ is finite
there exist finitely many ring maps
$b_1, \ldots, b_n : B \to k$ such that $b_i \circ s^\sharp = a_1$.
If the dotted arrow does not exist, then we see that none of the
$a'_i = b_i \circ t^\sharp$, $i = 1, \ldots, n$ is equal to $a_0$.
Hence the maximal ideals
$$
\mathfrak m'_i = \Ker(a_i' \otimes 1 : A \otimes_C k \to k)
$$
of $A \otimes_C k$ are distinct from
$\mathfrak m = \Ker(a_0 \otimes 1 : A \otimes_C k \to k)$.
By Algebra, Lemma \ref{algebra-lemma-silly} we would get an element
$f \in A \otimes_C k$ with $f \in \mathfrak m$, but
$f \not \in \mathfrak m_i'$ for $i = 1, \ldots, n$.
Consider the norm
$$
g = \text{Norm}_{s^\sharp \otimes 1}(t^\sharp \otimes 1(f))
\in
A \otimes_C k
$$
By Lemma \ref{lemma-determinant-trick} this lies in the invariants
$C^1 \subset A \otimes_C k$ of the base change
groupoid (base change via the map $c : C \to k$). On the one hand,
$a_1(g) \in k^*$ since
the value of $t^\sharp(f)$ at all the points (which correspond to
$b_1, \ldots, b_n$) lying over $a_1$ is
invertible (insert future reference on property determinant here).
On the other hand, since $f \in \mathfrak m$, we see that
$f$ is not a unit, hence $t^\sharp(f)$ is not a unit
(as $t^\sharp \otimes 1$ is faithfully flat),
hence its norm is not a unit (insert future reference
on property determinant here). We conclude that $C^1$ contains
an element which is not nilpotent
and not a unit. We will now show that this leads to a contradiction.
Namely, apply Lemma \ref{lemma-invariants-base-change}
to the map $c : C \to C' = k$, then
we see that the map of $k$ into the invariants $C^1$ is injective
and moreover, that for any element $x \in C^1$ there exists an integer
$n > 0$ such that $x^n \in k$. Hence every element of $C^1$ is
either a unit or nilpotent.

\medskip\noindent
We still have to finish the proof of (1). We already know that
$|U| \to |M|$ is surjective. It is clear that $|U| \to |M|$ is
$|R|$-invariant. Finally, suppose $u_0, u_1 \in U$ maps to the same
point $m \in M$. Then the induced field extensions $\kappa(u_0)/\kappa(m)$
and $\kappa(u_1)/\kappa(m)$ are algebraic (as $A$ is integral over $C$
as used above). Hence if $k$ is an algebraic closure of $\kappa(m)$
then we can find $\kappa(m)$-embeddings $\overline{u}_0 : \kappa(u_0) \to k$
and $\overline{u}_1 : \kappa(u_1) \to k$. These determine $k$-valued
points $\overline{u}_0, \overline{u}_1 \in U(k)$ mapping to the same
point of $M(k)$. By part (2) we see that there exists a point
$\overline{r} \in R(k)$ with $s(\overline{r}) = \overline{u}_0$ and
$t(\overline{r}) = \overline{u}_1$. The image $r \in R$ of $\overline{r}$
is a point with $s(r) = u_0$ and $t(r) = u_1$ as desired.
\end{proof}

\begin{proof}
Set $M = \Spec(C)$ and $M' = \Spec(C')$.
Write $U = \Spec(A)$, $U' = \Spec(A')$, $R = \Spec(B)$, and
$R' = \Spec(B')$. We will use the results of
Lemmas \ref{lemma-integral-over-invariants},
\ref{lemma-invariants-base-change}, and
\ref{lemma-points}
without further mention.

\medskip\noindent
Assume $C$ is a strictly henselian local ring. Let $p \in M$
be the closed point and let $p' \in M'$ map to $p$.
Claim: in this case there is a disjoint union decomposition
$(U', R', s', t', c') = (U, R, s, t, c) \amalg (U'', R'', s'', t'', c'')$
over $(U, R, s, t, c)$ such that for the corresponding
disjoint union decomposition $M' = M \amalg M''$ over $M$
the point $p'$ corresponds to $p \in M$.

\medskip\noindent
The claim implies the lemma. Suppose that $M_1 \to M$ is a flat morphism
of affine schemes. Then we can base change everything to $M_1$
without affecting the hypotheses (1) -- (4).
From Lemma \ref{lemma-invariants-base-change}
we see $M_1$, resp.\ $M_1'$ is the spectrum of the
$R_1$-invariant functions on $U_1$,
resp.\ the $R'_1$-invariant functions on $U'_1$.
Suppose that $p' \in M'$ maps to $p \in M$.
Let $M_1$ be the spectrum of the strict henselization of
$\mathcal{O}_{M, p}$ with closed point $p_1 \in M_1$.
Choose a point $p'_1 \in M'_1$ mapping to $p_1$ and $p'$.
From the claim we get
$$
(U'_1, R'_1, s'_1, t'_1, c'_1) =
(U_1, R_1, s_1, t_1, c_1) \amalg
(U''_1, R''_1, s''_1, t''_1, c''_1)
$$
and correspondingly $M'_1 = M_1 \amalg M''_1$ as a scheme over $M_1$.
Write $M_1 = \Spec(C_1)$ and write $C_1 = \colim C_i$ as a filtered
colimit of \'etale $C$-algebras. Set $M_i = \Spec(C_i)$.
The $M_1 = \lim M_i$ and similarly for the other schemes.
By Limits, Lemmas \ref{limits-lemma-descend-opens} and
\ref{limits-lemma-descend-isomorphism}
we can find an $i$ such that
$$
(U'_i, R'_i, s'_i, t'_i, c'_i) =
(U_i, R_i, s_i, t_i, c_i) \amalg
(U''_i, R''_i, s''_i, t''_i, c''_i)
$$
We conclude that $M'_i = M_i \amalg M''_i$. In particular
$M' \to M$ becomes \'etale at a point over $p'$ after an
\'etale base change. This implies that $M' \to M$ is \'etale at $p'$
(for example by Morphisms, Lemma
\ref{morphisms-lemma-set-points-where-fibres-etale}).
We will prove $U' \cong M' \times_M U$ after we prove the claim.

\medskip\noindent
Proof of the claim. Observe that $U_p$ and $U'_{p'}$ have finitely many points.
For $u \in U_p$ we have $\kappa(u)/\kappa(p)$ is algebraic,
hence $\kappa(u)$ is separably closed.
As $U' \to U$ is \'etale, we conclude the morphism $U'_{p'} \to U_p$
induces isomorphisms on residue field extensions.
Let $u' \in U'_{p'}$ with image $u \in U_p$.
By assumption (3) the morphism of scheme theoretic fibres
$(s')^{-1}(u') \to s^{-1}(u)$,
$(t')^{-1}(u') \to t^{-1}(u)$, and
$G'_{u'} \to G_u$ are isomorphisms. Observing that $U_p = t(s^{-1}(u))$
(set theoretically) we conclude that the points of $U'_{p'}$
surject onto the points of $U_p$.
Suppose that $u'_1$ and $u'_2$ are points of $U'_{p'}$ mapping
to the same point $u$ of $U_p$. Then there exists a point
$r' \in R'_{p'}$ with $s'(r') = u'_1$ and $t'(r') = u'_2$.
Consider the two towers of fields
$$
\kappa(r')/\kappa(u'_1)/\kappa(u)/\kappa(p) \quad
\kappa(r')/\kappa(u'_2)/\kappa(u)/\kappa(p)
$$
whose ``ends'' are the same as the two ``ends'' of the two towers
$$
\kappa(r')/\kappa(u'_1)/\kappa(p')/\kappa(p) \quad
\kappa(r')/\kappa(u'_2)/\kappa(p')/\kappa(p)
$$
These two induce the same maps $\kappa(p') \to \kappa(r')$ as
$(U'_{p'}, R'_{p'}, s', t', c')$ is a groupoid over $p'$.
Since $\kappa(u)/\kappa(p)$ is purely inseparable,
we conclude that the two induced maps
$\kappa(u) \to \kappa(r')$ are the same.
Therefore $r'$ maps to a point of the fibre $G_u$.
By assumption (3) we conclude that $r' \in (G')_{u'_1}$.
Namely, we may think of $G$ as a closed subscheme of $R$
viewed as a scheme over $U$ via $s$ and use that
the base change to $U'$ gives $G' \subset R'$.
In particular we have $u'_1 = u'_2$.
We conclude that $U'_{p'} \to U_p$ is a bijective
map on points inducing isomorphisms on residue fields.
It follows that $U'_{p'}$ is a finite set of closed points
(Algebra, Lemma \ref{algebra-lemma-finite-residue-extension-closed})
and hence $U'_{p'}$ is closed in $U'$.
Let $J' \subset A'$ be the radical ideal cutting out $U'_{p'}$
set theoretically.

\medskip\noindent
Second part proof of the claim.
Let $\mathfrak m \subset C$ be the maximal ideal.
Observe that $(A, \mathfrak m A)$ is a henselian pair by
More on Algebra, Lemma \ref{more-algebra-lemma-integral-over-henselian-pair}.
Let $J = \sqrt{\mathfrak m A}$.
Then $(A, J)$ is a henselian pair
(More on Algebra, Lemma \ref{more-algebra-lemma-change-ideal-henselian-pair})
and the \'etale ring map
$A \to A'$ induces an isomorphism $A/J \to A'/J'$
by our deliberations above.
We conclude that $A' = A \times A''$ by
More on Algebra, Lemma \ref{more-algebra-lemma-characterize-henselian-pair}.
Consider the corresponding disjoint union
decomposition $U' = U \amalg U''$. The open $(s')^{-1}(U)$ is the
set of points of $R'$ specializing to a point of $R'_{p'}$.
Similarly for $(t')^{-1}(U)$. Similarly we have
$(s')^{-1}(U'') = (t')^{-1}(U'')$ as this is the set of
points which do not specialize to $R'_{p'}$.
Hence we obtain a disjoint union decomposition
$$
(U', R', s', t', c') =
(U, R, s, t, c) \amalg
(U'', R'', s'', t'', c'')
$$
This immediately gives $M' = M \amalg M''$ and the proof of the claim
is complete.

\medskip\noindent
We still have to prove that the canonical map $U' \to M' \times_M U$
is an isomorphism. It is an \'etale morphism
(Morphisms, Lemma \ref{morphisms-lemma-etale-permanence}).
On the other hand, by base changing to strictly henselian local rings
(as in the third paragraph of the proof) and using the bijectivity
$U'_{p'} \to U_p$ established in the course of the proof of the claim,
we see that $U' \to M' \times_M U$ is universally bijective
(some details omitted). However, a universally bijective
\'etale morphism is an isomorphism
(Descent, Lemma \ref{descent-lemma-universally-injective-etale-open-immersion})
and the proof is complete.
\end{proof}

\bibliographystyle{plain}
\bibliography{my}
\end{document}
